% Sumber angka: results/eval_prediksi/gradient_boosting_h6.json, kunci folds dan
% rerata_lintas_fold. Jendela tersegmentasi (max_gap_steps=6). Versi tabel sebelumnya
% berasal dari jalan terdahulu yang hanya memuat dua model dan sudah digantikan.
%
% FORMAT (25 Agustus 2026): tabular -> tabularx selebar \textwidth. Sebelumnya keluar
% margin 55,4 pt. Dua sumber lebarnya dipangkas: awalan "ohio\_" pada kolom pasien uji
% dipindah ke keterangan tabel, dan satuan persen pada kepala kolom terakhir pun
% dipindah ke sana. Nilainya sendiri tidak berubah.
\begin{table}[htbp]
\centering
\caption{Validasi silang lintas-pasien (6 \emph{fold}, horizon $+30$ menit). Pada tiap \emph{fold}, ketiga model dilatih ulang dari nol pada 10 pasien dan diuji pada 2 pasien yang tak pernah dilihat.}
\label{tab:crossval}
\small
\begin{tabularx}{\textwidth}{@{}lLrrrr@{}}
\toprule
\textbf{Fold} & \textbf{Pasien uji} & \textbf{GBM} & \textbf{RF} & \textbf{LSTM} & \textbf{GBM A+B} \\
\midrule
0 & 540, 570 & 19,83 & 20,45 & 19,90 & 94,94 \\
1 & 544, 575 & 19,76 & 20,05 & 19,54 & 93,84 \\
2 & 552, 584 & 22,80 & 23,09 & 22,79 & 94,38 \\
3 & 559, 588 & 21,79 & 22,21 & 22,04 & 94,68 \\
4 & 563, 591 & 20,31 & 20,50 & 20,10 & 95,49 \\
5 & 567, 596 & 19,84 & 20,18 & 19,24 & 94,78 \\
\midrule
\textbf{Rerata} & -- & \textbf{20,72$\pm$1,16} & 21,08$\pm$1,15 & 20,60$\pm$1,33 & \textbf{94,69$\pm$0,50} \\
\bottomrule
\end{tabularx}

\vspace{0.4em}
{\tabnotestyle Kolom pasien uji memuat kode pasien OhioT1DM tanpa awalannya, sehingga ``540'' berarti \texttt{ohio\_\allowbreak 540}. Kolom GBM, RF, dan LSTM memuat RMSE dalam mg/dL, sedangkan kolom terakhir memuat proporsi zona aman Clarke~A+B dalam persen; baris terakhir memuat rerata beserta simpangan bakunya lintas \emph{fold}. \textbf{Selisih antar-model lebih kecil daripada simpangan baku antar-\emph{fold}} pada ketiga pasangan, sehingga menurut aturan pelaporan yang ditetapkan di muka pada penelitian ini, tidak satu pun di antaranya boleh dinarasikan sebagai keunggulan. Yang dapat disimpulkan adalah bahwa ketiganya \emph{setara} pada galat prediksi dan bahwa keragaman antar-pasien jauh melampaui keragaman antar-model. Dasar pemilihan model produksi karena itu tidak terletak pada tabel ini.\par}
\end{table}
